\section{Quantum Pseudocode: Additional Definitions}\label{sec:pseudocode-defs}
  Here, we show how operations on qubit registers are defined for
    our quantum pseudocode, as introduced in \cref{sec:pseudocode}.

  \subsection{Operating on Registers}
    First, we formally define for what it means for a unitary to be applied to
    a qubit register.
    The main idea is that, thanks to the requirement that the global state
    $\ket{\qcomp}$ is in a well-defined pure state, we can define an operation
    on sub-registers $\register$ of $\qcomp$ in terms of a larger operation on
    the whole state.

    \begin{definition}[Operating on registers]\label{def:apply}
      Let $\qcomp$ be an $n$-qubit quantum memory in quantum state
      $\ket{\qcomp}$, let the register $\register$ be an $m$-qubit
      subsystem of $\qcomp$, and let $\unitary$ be an
      $M \times M$ unitary matrix. 

      Next, define $\unitary_{\register}$ to be the $N \times N$ unitary matrix that
      has the effect of applying $\unitary$ on the qubits in $\register$ and the
      identity on the qubits in $\qcomp \backslash \register$.       
      Then, the application of $\unitary$ to the qubits in register $\register$ is
      captured in the method $\applyU{\unitary}{\register}[\qcomp]$, defined as
      \begin{align*}
        \applyU{\unitary}{\register}[\qcomp] &\coloneqq \ket{\qcomp(t)}
        \apply{\unitary_{\register}} \ket{\qcomp(t+1)}\, .
      \end{align*}
    \end{definition}

    \noindent
    Note how local operations affecting distant qubits through entanglement
      is captured by $\applyU{\unitary}{\register}[\qcomp]$ 
      being an evolution of the global (pure) state $\ket\qcomp$.
    
    To see that $\unitary_{\register}$ is well-defined, note that it may be
      implemented as follows:
      \[
        \unitary_{\register} \coloneqq \SORTgate^\dagger([\register], [\qcomp
        \backslash \register])\ (\unitary \otimes \idmatrix)\
        \SORTgate([\register], [\qcomp \backslash \register])\, ,
      \]
      where $\SORTgate([\register], [\qcomp \backslash \register])$ is the $N
      \times N$ unitary matrix, composed of a series of $\SWAPgate$ gates, that 
      reorders the qubits in $\qcomp$ such that the qubits of $\register$
      precede the qubits of $\qcomp \backslash \register$;
      $\SORTgate^\dagger$ reverses the operation of $\SORTgate$; and 
      $\idmatrix$ is the $2^{n-m} \times 2^{n-m}$ identity matrix.

  \subsection{Sending Registers}\label{subsec:sending}
    In order to say ``Alice sends a qubit register to Bob'', we 
      give Alice and Bob to the disjoint subregisters
      $\register[A]$ and $\register[B]$, respectively, 
      of the total system $\qcomp$. 
      Then, Alice sending register $\register[R]$ to Bob simply amounts to
      transferring $\register[R]$ from $\register[A]$ to $\register[B]$.

    \begin{definition}[Transferring a subregister]
      Let $\register[A]$ and $\register[B]$ be non-overlapping qubit registers. A
      subregister $\register[R]$ is said to be \emph{transferred} from
      $\register[A]$ to $\register[B]$, denoted $\register[A]
      \apply{\register[R]} \register[B]$, if the elements of $\register[R]$ are
      simultaneously removed from $\register[A]$ and added to $\register[B]$:
      \[
        \register[A] \apply{\register[R]} \register[B]
        \quad \Longleftrightarrow \quad 
        \register[A] \gets \register[A] \backslash \register[R]\, , \quad \register[B] \listapp \register[R]\, .
      \]
    \end{definition}

  \subsection{Measuring Registers}\label{sec:measurements}
    The final piece of the puzzle is defining the measurement operation
      $\meas{}{\register}$.

    \begin{definition}[Measuring a register]\label{def:measure}
      Let $\qcomp$ be an $n$-qubit quantum memory in pure state
        $\ket{\qcomp}$, let the register $\register$ be an $m$-qubit
        subsystem of $\qcomp$, and let $\measurement$ be an $m$-qubit
        projective measurement that check the state for membership in
        Hilbert subspaces $\{ \Sspace_i \}$.

      Next, define $\measurement_{\register}$ to be the $n$-qubit
        projective measurement that
        has the effect of applying $\measurement$ on the qubits in $\register$ and the
        identity on the qubits in $\qcomp \backslash \register$

      Then, performing measurement $\measurement$ on the qubits in register
        $\register$ is captured in the (probabilistic) method
        $\meas{\measurement}{\register}[\qcomp]$ which observes
        $\ket{\qcomp}$ under $\measurement_{\register}$. Observing outcome
        $i$, designated by the measurement procedure returning value $m_i$,
        is written
      \[
        i \sample \meas{\measurement}{\register}[\qcomp]\, ,
      \]
        and the post-measurement will be 
      \[
        \ket{\qcomp} \apply{\measurement} \ket{\qcomp}_{\Sspace_i}\, ,
      \]
      where $\ket{\qcomp}_{\Sspace_i}$ is the result of projecting the
        qubits of $\register$ onto subspace $\Sspace_i$, and applying the
        identity to the rest.
    \end{definition}

    \noindent
    Finally, observe that the projective measurement $\measurement_{\register}$ can be
      constructed similarly to $\unitary_{\register}$ above.
